\section{Ba lối thoát, và chữ ấy là chữ của sách}
\label{sec:ch16-ba-loi-thoat}

\index{lối thoát!chỗ mà cả ba đường thoát khỏi|(}%
Phần mở đầu chương vừa nói chữ \enquote{lối thoát} là chữ của sách, và nhận một
chữ về thì phải nói rõ nó chỉ cái gì. Ba đường ấy thoát khỏi đúng một chỗ, và
sau khi chỗ ấy được nêu thì hỏi được từng bài xem nó có tự đứng ở đó hay không.

\index{độ trung thực!chỗ mắc kẹt của định nghĩa}%
Chỗ cần thoát khỏi thì mục~\ref{sec:ch14-ground-truth-ke-tan-cong} đã phát biểu
gọn nhất. Định nghĩa~\ref{def:ch01-do-trung-thuc} hỏi lời giải thích có phản ánh
đúng những thứ mà hành vi của mô hình phụ thuộc vào hay không, và với một mô
hình có sẵn thì không ai biết những thứ ấy là gì. Vế trái của phép so sánh nằm
trong tay ta, vì lời giải thích là thứ ta vừa tính ra. Vế phải thì không. Mọi họ
\tn{chỉ số trung thực}{faithfulness metric} của
chương~\ref{ch:do-do-trung-thuc} tồn tại để né vế phải bằng một
\tn{đại lượng thay thế}{proxy}, và bảy chương của Phần~IV là bảy cách cho thấy
phép né ấy chưa ai kiểm.

Ba đường thoát khác nhau ở chỗ chúng làm gì với vế phải. Đường thứ nhất ép nó
phải biết trước: chọn hoặc dựng một thiết lập mà câu trả lời đúng là hệ quả của
chính thiết kế, rồi chỉ chấm điểm trong thiết lập ấy. Đường thứ hai thay hẳn vế
phải bằng một vế khác: đưa lời giải thích cho người đọc và hỏi họ. Đường thứ ba
đi thẳng vào vế phải: mở mô hình ra và đọc xem bên trong nó có gì. Sách đã gặp
cả ba, mỗi lần trong lúc đọc một bài báo khác nhau và cho một mục đích khác
nhau, nên đây là lần đầu ba đường đứng cạnh nhau.

\index{lối thoát!bài nào tự nhận là lối thoát}%
Bây giờ hỏi từng bài. Bài của chương~\ref{ch:chi-so-khong-do-do-trung-thuc} tự
đứng đúng ở chỗ ấy và nói ra: mục 3 của nó mở bằng đúng câu hỏi trên, rằng muốn
biết lời giải thích có trung thực không thì phải biết mô hình đã tính những gì,
rồi ghi rằng các tính toán bên trong không quan sát trực tiếp được, rồi đề xuất
thiết kế nhiệm vụ như cách vòng qua~\autocite{p21metrics}. Với bài ấy thì chữ
\enquote{lối thoát} là mô tả đúng và không cần ai gán cho nó.

Bài của chương~\ref{ch:con-nguoi-vang-mat} thì không. Nó không nhắc tới
\tn{độ trung thực}{faithfulness} lần nào, không nhắc ground truth lần nào,
và không nêu tên một chỉ số nào
của chương~\ref{ch:do-do-trung-thuc}. Cái nó đòi kiểm chứng là lời tuyên bố rằng
một phương pháp giải thích được cho người, tức phía
\tn{tính hợp lý}{plausibility} của phép chia mà
mục~\ref{sec:ch01-do-trung-thuc-va-tinh-hop-ly} dựng. Nó là lối thoát khỏi một
chỗ mắc kẹt khác, và mục~\ref{sec:ch16-duong-thu-hai} sẽ cho thấy cách
đọc lẫn hai chỗ ấy làm hỏng điều gì.

Hai bài của đường thứ ba thì không tự nhận là lối thoát khỏi chỗ nào cả. Bài của
mục~\ref{sec:ch16-duong-thu-ba} là một bài quan điểm về attribution, xếp
\tn{diễn giải cơ chế}{mechanistic interpretability} vào bản đồ chứ không làm nó;
bài còn lại đề xuất một cách sửa cho chỗ yếu của chính loại công cụ mà nó dùng,
tức sparse autoencoder mà mục~\ref{sec:ch15-doc-duoc-dieu-khien-duoc} đã đọc.
Chữ \enquote{faithful} không xuất hiện lần nào trong phần thân của cả hai.

\index{lối thoát!ba đường và chữ của sách}%
Vậy trong bốn bài thì một bài đứng đúng chỗ nhan đề chương đặt nó, một bài đứng
ở chỗ khác, và hai bài không đứng ở chỗ nào. Sách vẫn giữ cách sắp xếp ba đường,
vì nó là cách sắp xếp của chính văn liệu phê phán mà
mục~\ref{sec:ch01-luan-de-cua-sach} đã nêu khi phát biểu luận đề, nhưng giữ nó
như cách sắp xếp của sách. Không câu nào trong sáu mục sau được viết như thể một
bài báo tự gọi mình là lối thoát.

\index{lối thoát!cái giá nghĩa là gì ở đây}%
Còn một chữ nữa phải định nghĩa trước khi dùng, và nó là chữ mà chương này sửa
lại. \enquote{Cái giá} của một lối thoát, ở đây, không phải chi phí tính toán và
cũng không phải tiền. Nó là câu trả lời cho một câu hỏi: sau khi đi đường ấy thì
ta biết được gì, và cái ta biết có phải cái định nghĩa~\ref{def:ch01-do-trung-thuc}
hỏi hay không. Ba đường cho ba câu trả lời khác nhau về hạng, và
mục~\ref{sec:ch16-ba-cai-gia} đặt cả ba cạnh nhau sau khi đã đọc từng đường.

\index{lối thoát!chỗ mà cả ba đường thoát khỏi|)}%
Cách làm của chương cũng khác các chương trước và đáng nói một câu, vì nó giải
thích vì sao một chương không đọc bài mới nào lại tìm ra được gì. Bốn bài ở đây
đều đã được đọc kỹ: bài của đường thứ nhất ở
chương~\ref{ch:do-do-trung-thuc} và chương~\ref{ch:chi-so-khong-do-do-trung-thuc},
bài của đường thứ hai ở chương~\ref{ch:con-nguoi-vang-mat}, và hai bài của đường
thứ ba ở chương~\ref{ch:khung-hop-nhat} và chương~\ref{ch:nut-that-khai-niem}.
Nhưng mỗi chương ấy đọc bài của mình cho một câu hỏi riêng, và không chương nào
hỏi câu của chương này. Đọc lại bốn bài cho đúng một câu hỏi mới là việc mà mục
sau bắt đầu.
